\balance
\begin{thebibliography}{99}\small
\bibitem{Chalmers1995} D. J. Chalmers, ``Facing up to the problem of consciousness,'' \emph{J. Conscious. Stud.}, vol. 2, no. 3, pp. 200--219, 1995.
\bibitem{Rosenthal1993} D. M. Rosenthal, ``State consciousness and transitive consciousness,'' \emph{Conscious. Cogn.}, vol. 2, no. 4, pp. 355--363, 1993, doi:10.1006/ccog.1993.1029.
\bibitem{Block1995} N. Block, ``On a confusion about a function of consciousness,'' \emph{Behav. Brain Sci.}, vol. 18, no. 2, pp. 227--247, 1995, doi:10.1017/S0140525X00038188.
\bibitem{Sanders2012} R. D. Sanders, G. Tononi, S. Laureys, and J. W. Sleigh, ``Unresponsiveness $\neq$ unconsciousness,'' \emph{Anesthesiology}, vol. 116, no. 4, pp. 946--959, 2012, doi:10.1097/ALN.0b013e318249d0a7.
\bibitem{Bayne2016} T. Bayne, J. Hohwy, and A. M. Owen, ``Are there levels of consciousness?'' \emph{Trends Cogn. Sci.}, vol. 20, no. 6, pp. 405--413, 2016, doi:10.1016/j.tics.2016.03.009.
\bibitem{LeDoux2025} J. E. LeDoux, ``What the functions of consciousness are depends on what one thinks consciousness is,'' \emph{Philos. Trans. R. Soc. B}, vol. 380, art. 20240311, 2025, doi:10.1098/rstb.2024.0311.
\bibitem{Noreika2011} V. Noreika et al., ``Consciousness lost and found: subjective experiences in an unresponsive state,'' \emph{Brain Cogn.}, vol. 77, no. 3, pp. 327--334, 2011, doi:10.1016/j.bandc.2011.09.002.
\bibitem{Casey2024} C. P. Casey et al., ``Evaluation of putative signatures of consciousness using specific definitions of responsiveness, connectedness, and consciousness,'' \emph{Br. J. Anaesth.}, vol. 132, no. 2, pp. 300--311, 2024, doi:10.1016/j.bja.2023.09.031.
\bibitem{Bodien2024} Y. G. Bodien et al., ``Cognitive motor dissociation in disorders of consciousness,'' \emph{N. Engl. J. Med.}, vol. 391, pp. 598--608, 2024, doi:10.1056/NEJMoa2400645.
\bibitem{Milner2005} B. Milner, ``The medial temporal-lobe amnesic syndrome,'' \emph{Psychiatr. Clin. North Am.}, vol. 28, no. 3, pp. 599--611, 2005, doi:10.1016/j.psc.2005.06.002.
\bibitem{Fulda2023} F. C. Fulda, ``Agential autonomy and biological individuality,'' \emph{Evol. Dev.}, vol. 25, no. 6, pp. 353--370, 2023, doi:10.1111/ede.12450.
\bibitem{BechtelBich2024} W. Bechtel and L. Bich, ``Situating homeostasis in organisms: maintaining organization through time,'' \emph{J. Physiol.}, vol. 602, no. 22, pp. 6003--6020, 2024, doi:10.1113/JP286883.
\bibitem{Bronfman2016} Z. Z. Bronfman, S. Ginsburg, and E. Jablonka, ``The transition to minimal consciousness through the evolution of associative learning,'' \emph{Front. Psychol.}, vol. 7, art. 1954, 2016, doi:10.3389/fpsyg.2016.01954.
\bibitem{Birch2020} J. Birch, S. Ginsburg, and E. Jablonka, ``Unlimited associative learning and the origins of consciousness: a primer and some predictions,'' \emph{Biol. Philos.}, vol. 35, art. 56, 2020, doi:10.1007/s10539-020-09772-0.
\bibitem{Greer2023} D. M. Greer et al., ``Pediatric and adult brain death/death by neurologic criteria consensus guideline,'' \emph{Neurology}, vol. 101, no. 24, pp. 1112--1132, 2023, doi:10.1212/WNL.0000000000207740.
\bibitem{Shewmon1998} D. A. Shewmon, ``Chronic `brain death': meta-analysis and conceptual consequences,'' \emph{Neurology}, vol. 51, no. 6, pp. 1538--1545, 1998, doi:10.1212/WNL.51.6.1538.
\end{thebibliography}

\vspace{0.5em}
\noindent\footnotesize\textbf{Methodological note.} This article is a theoretical and methodological proposal. It reports no new human or animal experiment and is not clinical, legal, or ethical guidance. Empirical claims are limited to the cited literature. The definition in Eq. \ref{eq:def} is revisionary and is evaluated by coherence, comparative scientific utility, and the explicit failure conditions stated above; it is not presented as a result established by the dissociation cases themselves.\\[0.35em]\textbf{Open materials.} Manuscript PDF, LaTeX source, experiment specifications, reviewer-response notes, social-post copy, and explainer-video materials are archived at \href{\repoURL}{the project repository}.

\end{document}
